\import{nat/multiplication.tex}

\section{Natural divisibility and greatest common divisors}

The arithmetic in this section is wholly internal to the natural numbers.
In particular, the proofs use neither signed coefficients nor the surreal
integer implementation.

\begin{definition}\label{natural_one}
    $\naturalOne=\suc{\emptyset}$.
\end{definition}

\begin{definition}\label{natural_divisibility}
    $m$ divides $n$ as natural numbers iff
        $m$ is a natural number and $n$ is a natural number and
        there exists a natural number $k$ such that
            $n=\naturalMul{m}{k}$.
\end{definition}

\begin{proposition}\label{natural_divisibility_guards}
    Suppose $m$ divides $n$ as natural numbers.
    Then $m$ is a natural number and $n$ is a natural number.
\end{proposition}
\begin{proof}
    Follows by \obv{natural_divisibility}.
\end{proof}

\begin{definition}\label{common_divisors}
    $\naturalCommonDivisors{m}{n}=
        \{d\in\suc{m\union n}\mid
            \text{$m$ and $n$ are natural numbers and
                $d$ divides $m$ as natural numbers and
                $d$ divides $n$ as natural numbers}\}$.
\end{definition}

\begin{definition}\label{natural_gcd}
    $\naturalGcd{m}{n}=\unions{\naturalCommonDivisors{m}{n}}$.
\end{definition}

\begin{definition}\label{coprime_naturals}
    $m$ is coprime to $n$ as natural numbers iff
        $m$ is a natural number and $n$ is a natural number and
        $\naturalGcd{m}{n}=\naturalOne$.
\end{definition}

\begin{proposition}\label{coprime_naturals_guards}
    Suppose $m$ is coprime to $n$ as natural numbers.
    Then $m$ is a natural number and $n$ is a natural number.
\end{proposition}
\begin{proof}
    Follows by \obv{coprime_naturals}.
\end{proof}

\begin{definition}\label{prime_natural}
    $p$ is a prime natural number iff
        $p$ is a natural number and $\naturalOne\in p$ and
        for every natural number $d$ such that
            $d$ divides $p$ as natural numbers we have
            $d=\naturalOne$ or $d=p$.
\end{definition}

\begin{proposition}\label{prime_natural_guard}
    Suppose $p$ is a prime natural number.
    Then $p$ is a natural number.
\end{proposition}
\begin{proof}
    Follows by \obv{prime_natural}.
\end{proof}

\subsection{Elementary divisibility calculus}

\begin{proposition}\label{natural_divides_itself}
    Let $m$ be a natural number.
    Then $m$ divides $m$ as natural numbers.
\end{proposition}
\begin{proof}
    $\suc{\emptyset}$ is a natural number
        by \obv{omega_empty,omega_successor}.
    $m=\naturalMul{m}{\suc{\emptyset}}$
        by \obv{natural_multiplication_one_right}.
    Follows by \obv{natural_divisibility}.
\end{proof}

\begin{proposition}\label{every_natural_divides_zero}
    Let $m$ be a natural number.
    Then $m$ divides $\emptyset$ as natural numbers.
\end{proposition}
\begin{proof}
    $\emptyset$ is a natural number by \obv{omega_empty}.
    $\naturalMul{m}{\emptyset}=\emptyset$
        by \obv{natural_multiplication_zero_right}.
    Follows by \obv{natural_divisibility}.
\end{proof}

\begin{proposition}\label{one_divides_every_natural}
    Let $n$ be a natural number.
    Then $\naturalOne$ divides $n$ as natural numbers.
\end{proposition}
\begin{proof}
    $\naturalOne$ is a natural number
        by \auto{natural_one,omega_empty,omega_successor}.
    $\naturalMul{\naturalOne}{n}=n$
        by \auto{natural_one,
            natural_multiplication_one_left}.
    Follows by \obv{natural_divisibility}.
\end{proof}

\begin{proposition}\label{zero_divides_natural_iff_zero}
    Let $n$ be a natural number.
    Then $\emptyset$ divides $n$ as natural numbers iff
        $n=\emptyset$.
\end{proposition}
\begin{proof}
    Show that if $\emptyset$ divides $n$ as natural numbers, then
        $n=\emptyset$.
    \begin{subproof}
        Assume $\emptyset$ divides $n$ as natural numbers.
        Take a natural number $k$ such that
            $n=\naturalMul{\emptyset}{k}$
            by \obv{natural_divisibility}.
        Follows by \obv{natural_multiplication_zero_left}.
    \end{subproof}
    Show that if $n=\emptyset$, then
        $\emptyset$ divides $n$ as natural numbers.
    \begin{subproof}
        Assume $n=\emptyset$.
        Follows by \obv{every_natural_divides_zero,omega_empty}.
    \end{subproof}
    Follows by \auto.
\end{proof}

\begin{proposition}\label{natural_divisibility_transitive}
    Let $d,m,n$ be natural numbers.
    Suppose $d$ divides $m$ as natural numbers and
        $m$ divides $n$ as natural numbers.
    Then $d$ divides $n$ as natural numbers.
\end{proposition}
\begin{proof}
    Take $a,b\in\naturals$ such that
        $m=\naturalMul{d}{a}$ and $n=\naturalMul{m}{b}$
        by \auto{natural_divisibility}.
    $\naturalMul{a}{b}$ is a natural number
        by \obv{natural_multiplication_closed}.
    $n=\naturalMul{d}{\naturalMul{a}{b}}$
        by \auto{natural_multiplication_associative}.
    Follows by \obv{natural_divisibility}.
\end{proof}

\begin{proposition}\label{common_divisor_divides_sum}
    Let $d,m,n$ be natural numbers.
    Suppose $d$ divides $m$ as natural numbers and
        $d$ divides $n$ as natural numbers.
    Then $d$ divides $m+n$ as natural numbers.
\end{proposition}
\begin{proof}
    Take $a,b\in\naturals$ such that
        $m=\naturalMul{d}{a}$ and $n=\naturalMul{d}{b}$
        by \auto{natural_divisibility}.
    $a+b$ is a natural number by \obv{num_addition_on_naturals}.
    $m+n$ is a natural number by \obv{num_addition_on_naturals}.
    $m+n=\naturalMul{d}{a+b}$
        by \auto{natural_multiplication_distributive_left}.
    Follows by \obv{natural_divisibility}.
\end{proof}

\begin{proposition}\label{natural_multiplication_zero_iff}
    Let $a,b$ be natural numbers.
    Then $\naturalMul{a}{b}=\emptyset$ iff
        $a=\emptyset$ or $b=\emptyset$.
\end{proposition}
\begin{proof}
    Show that if $\naturalMul{a}{b}=\emptyset$, then
        $a=\emptyset$ or $b=\emptyset$.
    \begin{subproof}
        Assume $\naturalMul{a}{b}=\emptyset$.
        \begin{byCase}
            \caseOf{$b=\emptyset$.}
                Follows.
            \caseOf{$b\neq\emptyset$.}
                Take a natural number $c$ such that $b=\suc{c}$
                    by \obv{natural_zero_or_successor}.
                $\naturalMul{a}{b}=\naturalMul{a}{c}+a$
                    by \auto{natural_multiplication_successor_right}.
                $\naturalMul{a}{c}$ is a natural number
                    by \obv{natural_multiplication_closed}.
                $a\subseteq a+\naturalMul{a}{c}$
                    by \obv{natural_addition_left_subseteq_sum}.
                $a+\naturalMul{a}{c}=\naturalMul{a}{c}+a$
                    by \obv{natural_addition_commutative}.
                Hence $a\subseteq\naturalMul{a}{b}$.
                Thus $a\subseteq\emptyset$.
                Hence $a=\emptyset$
                    by \obv{subseteq_emptyset_iff}.
        \end{byCase}
    \end{subproof}
    Show that if $a=\emptyset$ or $b=\emptyset$, then
        $\naturalMul{a}{b}=\emptyset$.
    \begin{subproof}
        Assume $a=\emptyset$ or $b=\emptyset$.
        Follows by
            \auto{natural_multiplication_zero_left,
                natural_multiplication_zero_right}.
    \end{subproof}
    Follows by \auto.
\end{proof}

\begin{proposition}\label{positive_summand_makes_strict_sum}
    Let $a,b$ be natural numbers.
    Suppose $a\neq\emptyset$.
    Then $b\in a+b$.
\end{proposition}
\begin{proof}
    Take a natural number $c$ such that $a=\suc{c}$
        by \obv{natural_zero_or_successor}.
    $c+b$ is a natural number by \obv{num_addition_on_naturals}.
    $b\subseteq b+c$ by \obv{natural_addition_left_subseteq_sum}.
    Hence $b\subseteq c+b$ by \obv{natural_addition_commutative}.
    Hence $b\in c+b$ or $b=c+b$
        by \obv{natural_number_is_ordinal,
            ordinal_subseteq_implies_elem_or_eq}.
    Thus $b\in\suc{c+b}$ by \auto{suc_iff}.
    $a+b=\suc{c+b}$ by \auto{num_addition_successor_left}.
\end{proof}

\begin{theorem}\label{natural_multiplication_cancel_left}
    Let $d,a,b$ be natural numbers.
    Suppose $d\neq\emptyset$ and
        $\naturalMul{d}{a}=\naturalMul{d}{b}$.
    Then $a=b$.
\end{theorem}
\begin{proof}
    Let $C=\{x\in\naturals\mid
        \text{for every $y\in\naturals$ such that
            $\naturalMul{d}{x}=\naturalMul{d}{y}$ we have $x=y$}\}$.
    Show that $\emptyset\in C$.
    \begin{subproof}
        Show that for every $y\in\naturals$ such that
            $\naturalMul{d}{\emptyset}=\naturalMul{d}{y}$
            we have $\emptyset=y$.
        \begin{subproof}
            Fix $y\in\naturals$.
            Assume $\naturalMul{d}{\emptyset}=\naturalMul{d}{y}$.
            Hence $\naturalMul{d}{y}=\emptyset$
                by \auto{natural_multiplication_zero_right}.
            Hence $d=\emptyset$ or $y=\emptyset$
                by \obv{natural_multiplication_zero_iff}.
        \end{subproof}
        Follows by \obv{omega_empty}.
    \end{subproof}
    Show that for every $x\in C$ we have $\suc{x}\in C$.
    \begin{subproof}
        Fix $x\in C$.
        $\suc{x}$ is a natural number by \obv{omega_successor}.
        Show that for every $y\in\naturals$ such that
            $\naturalMul{d}{\suc{x}}=\naturalMul{d}{y}$
            we have $\suc{x}=y$.
        \begin{subproof}
            Fix $y\in\naturals$.
            Assume $\naturalMul{d}{\suc{x}}=\naturalMul{d}{y}$.
            Show that $y\neq\emptyset$.
            \begin{subproof}
                Suppose $y=\emptyset$.
                Then $\naturalMul{d}{\suc{x}}=\emptyset$
                    by \obv{natural_multiplication_zero_right}.
                Hence $d=\emptyset$ or $\suc{x}=\emptyset$
                    by \obv{natural_multiplication_zero_iff}.
                Contradiction by \obv{suc_neq_emptyset}.
            \end{subproof}
            Take a natural number $z$ such that $y=\suc{z}$
                by \obv{natural_zero_or_successor}.
            $\naturalMul{d}{x}+d=\naturalMul{d}{z}+d$
                by \auto{natural_multiplication_successor_right}.
            Hence $\naturalMul{d}{x}=\naturalMul{d}{z}$
                by \obv{natural_addition_cancel_right,
                    natural_multiplication_closed}.
            Hence $x=z$.
        \end{subproof}
    \end{subproof}
    Hence $C$ is an inductive set by \obv{inductive_set}.
    Hence every natural number is an element of $C$
        by \obv{omega_minimal}.
    Thus $a\in C$.
\end{proof}

\begin{theorem}\label{positive_multiplication_reflects_order}
    Let $d,a,b$ be natural numbers.
    Suppose $d\neq\emptyset$ and
        $\naturalMul{d}{a}\subseteq\naturalMul{d}{b}$.
    Then $a\subseteq b$.
\end{theorem}
\begin{proof}
    $a$ and $b$ are ordinals by \obv{natural_number_is_ordinal}.
    \begin{byCase}
        \caseOf{$a\subseteq b$.}
            Follows.
        \caseOf{$\lnot(a\subseteq b)$.}
            Hence $b\subseteq a$
                by \obv{ordinal_subseteq_or_subseteq}.
            Take a natural number $k$ such that $a=b+k$
                by \obv{natural_additive_difference}.
            $\naturalMul{d}{a}
                =\naturalMul{d}{b}+\naturalMul{d}{k}$
                by \auto{natural_multiplication_distributive_left}.
            $\naturalMul{d}{b}\subseteq\naturalMul{d}{a}$
                by \auto{natural_addition_left_subseteq_sum,
                    natural_multiplication_closed}.
            Hence $\naturalMul{d}{a}=\naturalMul{d}{b}$
                by \obv{subseteq_antisymmetric}.
            Hence $a=b$
                by \obv{natural_multiplication_cancel_left}.
    \end{byCase}
\end{proof}

\begin{proposition}\label{divisor_bounded_by_nonzero_dividend}
    Let $d,n$ be natural numbers.
    Suppose $n\neq\emptyset$ and
        $d$ divides $n$ as natural numbers.
    Then $d\subseteq n$.
\end{proposition}
\begin{proof}
    Take a natural number $k$ such that $n=\naturalMul{d}{k}$
        by \obv{natural_divisibility}.
    Show that $k\neq\emptyset$.
    \begin{subproof}
        Suppose $k=\emptyset$.
        Then $n=\emptyset$
            by \auto{natural_multiplication_zero_right}.
        Contradiction.
    \end{subproof}
    Take a natural number $l$ such that $k=\suc{l}$
        by \obv{natural_zero_or_successor}.
    $n=\naturalMul{d}{l}+d$
        by \auto{natural_multiplication_successor_right}.
    $\naturalMul{d}{l}$ is a natural number
        by \obv{natural_multiplication_closed}.
    $d\subseteq d+\naturalMul{d}{l}$
        by \obv{natural_addition_left_subseteq_sum}.
    $d+\naturalMul{d}{l}=\naturalMul{d}{l}+d$
        by \obv{natural_addition_commutative}.
\end{proof}

\begin{proposition}\label{larger_natural_does_not_divide_positive_smaller}
    Let $k,d$ be natural numbers.
    Suppose $k\neq\emptyset$ and $k\in d$.
    Then $d$ does not divide $k$ as natural numbers.
\end{proposition}
\begin{proof}
    Suppose that $d$ divides $k$ as natural numbers.
    Then $d\subseteq k$ by
        \obv{divisor_bounded_by_nonzero_dividend}.
    Hence $k\in k$ by \obv{subseteq}.
    Contradiction by \obv{in_irrefl}.
\end{proof}

\begin{theorem}\label{common_divisor_of_sum_difference}
    Let $d,m,k$ be natural numbers.
    Suppose $d$ divides $m$ as natural numbers and
        $d$ divides $m+k$ as natural numbers.
    Then $d$ divides $k$ as natural numbers.
\end{theorem}
\begin{proof}
    Take $a,b\in\naturals$ such that
        $m=\naturalMul{d}{a}$ and
        $m+k=\naturalMul{d}{b}$
        by \auto{natural_divisibility}.
    \begin{byCase}
        \caseOf{$d=\emptyset$.}
            Then $m=\emptyset$ and $m+k=\emptyset$
                by \auto{natural_multiplication_zero_left}.
            Hence $k=\emptyset$ by \auto{num_addition_zero_left}.
            Follows by \obv{every_natural_divides_zero,omega_empty}.
        \caseOf{$d\neq\emptyset$.}
            $m\subseteq m+k$ by
                \obv{natural_addition_left_subseteq_sum}.
            Hence $\naturalMul{d}{a}\subseteq\naturalMul{d}{b}$.
            Thus $a\subseteq b$
                by \obv{positive_multiplication_reflects_order}.
            Take a natural number $c$ such that $b=a+c$
                by \obv{natural_additive_difference}.
            \begin{align*}
                m+k
                    &= \naturalMul{d}{a}+\naturalMul{d}{c}
                        \explanation{by
                            \auto{natural_multiplication_distributive_left}}
                    \\
                    &= m+\naturalMul{d}{c}
            \end{align*}
            Thus $k=\naturalMul{d}{c}$
                by \auto{natural_addition_cancel_left,
                    natural_multiplication_closed}.
            Follows by \obv{natural_divisibility}.
    \end{byCase}
\end{proof}

\begin{corollary}\label{common_divisor_sum_iff_difference}
    Let $d,m,k$ be natural numbers.
    Then
        ($d$ divides $m$ as natural numbers and
            $d$ divides $m+k$ as natural numbers)
        iff
        ($d$ divides $m$ as natural numbers and
            $d$ divides $k$ as natural numbers).
\end{corollary}
\begin{proof}
    Show that if $d$ divides $m$ as natural numbers and
        $d$ divides $m+k$ as natural numbers,
        then $d$ divides $m$ as natural numbers and
            $d$ divides $k$ as natural numbers.
    \begin{subproof}
        Assume $d$ divides $m$ as natural numbers and
            $d$ divides $m+k$ as natural numbers.
        Follows by \obv{common_divisor_of_sum_difference}.
    \end{subproof}
    Show that if $d$ divides $m$ as natural numbers and
        $d$ divides $k$ as natural numbers,
        then $d$ divides $m$ as natural numbers and
            $d$ divides $m+k$ as natural numbers.
    \begin{subproof}
        Assume $d$ divides $m$ as natural numbers and
            $d$ divides $k$ as natural numbers.
        Follows by \obv{common_divisor_divides_sum}.
    \end{subproof}
    Follows by \auto.
\end{proof}

\subsection{The bounded greatest common divisor}

\begin{proposition}\label{common_divisors_iff}
    $d\in\naturalCommonDivisors{m}{n}$ iff
        $d\in\suc{m\union n}$ and
        $m$ is a natural number and $n$ is a natural number and
        $d$ divides $m$ as natural numbers and
        $d$ divides $n$ as natural numbers.
\end{proposition}
\begin{proof}
    Follows by \auto{common_divisors}.
\end{proof}

\begin{proposition}\label{gcd_is_union_of_common_divisors}
    $\naturalGcd{m}{n}=
        \unions{\naturalCommonDivisors{m}{n}}$.
\end{proposition}
\begin{proof}
    Follows by \obv{natural_gcd}.
\end{proof}

\begin{proposition}\label{common_divisors_symmetric}
    $\naturalCommonDivisors{m}{n}
        =\naturalCommonDivisors{n}{m}$.
\end{proposition}
\begin{proof}
    $m\union n=n\union m$ by \obv{union_comm}.
    Follows by
        \obv{common_divisors_iff,setext}.
\end{proof}

\begin{proposition}\label{natural_gcd_symmetric}
    $\naturalGcd{m}{n}=\naturalGcd{n}{m}$.
\end{proposition}
\begin{proof}
    Follows by
        \obv{gcd_is_union_of_common_divisors,
            common_divisors_symmetric}.
\end{proof}

\begin{proposition}\label{common_divisor_in_bounded_set}
    Let $d,m,n$ be natural numbers.
    Suppose $m\neq\emptyset$ and
        $d$ divides $m$ as natural numbers and
        $d$ divides $n$ as natural numbers.
    Then $d\in\naturalCommonDivisors{m}{n}$.
\end{proposition}
\begin{proof}
    $d\subseteq m$
        by \obv{divisor_bounded_by_nonzero_dividend}.
    $m\subseteq m\union n$ by \obv{union_upper_left}.
    Hence $d\subseteq m\union n$ by \obv{subseteq_transitive}.
    $d$ and $m\union n$ are ordinals
        by \obv{natural_number_is_ordinal,
            union_of_two_ordinals_is_ordinal}.
    Hence $d\in m\union n$ or $d=m\union n$
        by \obv{ordinal_subseteq_implies_elem_or_eq}.
    Thus $d\in\suc{m\union n}$ by \auto{suc_iff}.
    Follows by \obv{common_divisors_iff}.
\end{proof}

\begin{proposition}\label{natural_gcd_zero_zero}
    $\naturalGcd{\emptyset}{\emptyset}=\emptyset$.
\end{proposition}
\begin{proof}
    Show that $\naturalGcd{\emptyset}{\emptyset}\subseteq\emptyset$.
    \begin{subproof}
        Follows by
            \auto{gcd_is_union_of_common_divisors,
                common_divisors_iff,subseteq,unions_iff,
                union_emptyset,suc_iff,emptyset}.
    \end{subproof}
    Follows by \obv{subseteq_emptyset_iff}.
\end{proof}

\begin{proposition}\label{natural_gcd_positive_zero}
    Let $m$ be a natural number.
    Suppose $m\neq\emptyset$.
    Then $\naturalGcd{m}{\emptyset}=m$.
\end{proposition}
\begin{proof}
    Let $C=\naturalCommonDivisors{m}{\emptyset}$.
    Show that for every $x\in\unions{C}$ we have $x\in m$.
    \begin{subproof}
        Fix $x\in\unions{C}$.
        Take $d\in C$ such that $x\in d$ by \obv{unions_iff}.
        $d$ is a natural number and
            $d$ divides $m$ as natural numbers
            by \auto{common_divisors_iff,
                natural_divisibility}.
        Hence $d\subseteq m$
            by \obv{divisor_bounded_by_nonzero_dividend}.
        Follows by \obv{subseteq}.
    \end{subproof}
    Hence $\unions{C}\subseteq m$ by \obv{subseteq}.
    $m$ divides $m$ as natural numbers
        by \obv{natural_divides_itself}.
    $m$ divides $\emptyset$ as natural numbers
        by \obv{every_natural_divides_zero}.
    Hence $m\in C$
        by \obv{common_divisor_in_bounded_set}.
    Thus $m\subseteq\unions{C}$ by \auto{subseteq,unions_iff}.
    Hence $\unions{C}=m$ by \obv{subseteq_antisymmetric}.
    Follows by
        \obv{gcd_is_union_of_common_divisors}.
\end{proof}

\begin{theorem}\label{natural_gcd_zero_right}
    Let $m$ be a natural number.
    Then $\naturalGcd{m}{\emptyset}=m$.
\end{theorem}
\begin{proof}
    \begin{byCase}
        \caseOf{$m=\emptyset$.}
            Follows by \obv{natural_gcd_zero_zero}.
        \caseOf{$m\neq\emptyset$.}
            Follows by \obv{natural_gcd_positive_zero}.
    \end{byCase}
\end{proof}

\begin{corollary}\label{natural_gcd_zero_left}
    Let $n$ be a natural number.
    Then $\naturalGcd{\emptyset}{n}=n$.
\end{corollary}
\begin{proof}
    Follows by
        \auto{natural_gcd_symmetric,
            natural_gcd_zero_right}.
\end{proof}

\begin{theorem}\label{common_divisors_subtract_sum}
    Let $m,k$ be natural numbers.
    Suppose $m\neq\emptyset$.
    Then $\naturalCommonDivisors{m}{m+k}
        =\naturalCommonDivisors{m}{k}$.
\end{theorem}
\begin{proof}
    Show that for every
        $d\in\naturalCommonDivisors{m}{m+k}$ we have
        $d\in\naturalCommonDivisors{m}{k}$.
    \begin{subproof}
        Fix $d\in\naturalCommonDivisors{m}{m+k}$.
        Then $d$ is a natural number and
            $d$ divides $m$ as natural numbers and
            $d$ divides $m+k$ as natural numbers
            by \auto{common_divisors_iff,
                natural_divisibility}.
        Hence $d$ divides $k$ as natural numbers
            by \obv{common_divisor_of_sum_difference}.
        Follows by
            \obv{common_divisor_in_bounded_set}.
    \end{subproof}
    Hence $\naturalCommonDivisors{m}{m+k}
        \subseteq\naturalCommonDivisors{m}{k}$
        by \obv{subseteq}.
    Show that for every
        $d\in\naturalCommonDivisors{m}{k}$ we have
        $d\in\naturalCommonDivisors{m}{m+k}$.
    \begin{subproof}
        Fix $d\in\naturalCommonDivisors{m}{k}$.
        Then $d$ is a natural number and
            $d$ divides $m$ as natural numbers and
            $d$ divides $k$ as natural numbers
            by \auto{common_divisors_iff,
                natural_divisibility}.
        Hence $d$ divides $m+k$ as natural numbers
            by \obv{common_divisor_divides_sum}.
        Follows by
            \obv{common_divisor_in_bounded_set}.
    \end{subproof}
    Hence $\naturalCommonDivisors{m}{k}
        \subseteq\naturalCommonDivisors{m}{m+k}$
        by \obv{subseteq}.
    Follows by \obv{subseteq_antisymmetric}.
\end{proof}

\begin{corollary}\label{natural_gcd_subtract_sum}
    Let $m,k$ be natural numbers.
    Suppose $m\neq\emptyset$.
    Then $\naturalGcd{m}{m+k}=\naturalGcd{m}{k}$.
\end{corollary}
\begin{proof}
    Follows by
        \obv{gcd_is_union_of_common_divisors,
            common_divisors_subtract_sum}.
\end{proof}

\subsection{The divisibility universal property}

\begin{definition}\label{gcd_good_pairs}
    $\naturalGcdGoodPairs=
        \{p\in\naturals\times\naturals\mid
            \text{$\naturalGcd{\fst{p}}{\snd{p}}$ divides $\fst{p}$
                as natural numbers and
                $\naturalGcd{\fst{p}}{\snd{p}}$ divides $\snd{p}$
                as natural numbers and
                for every $d\in\naturals$ we have
                ($d$ divides $\naturalGcd{\fst{p}}{\snd{p}}$
                    as natural numbers iff
                 $d$ divides $\fst{p}$ as natural numbers and
                 $d$ divides $\snd{p}$ as natural numbers)}\}$.
\end{definition}

\begin{proposition}\label{gcd_good_pair_iff}
    $(m,n)\in\naturalGcdGoodPairs$ iff
        $m$ is a natural number and $n$ is a natural number and
        $\naturalGcd{m}{n}$ divides $m$ as natural numbers and
        $\naturalGcd{m}{n}$ divides $n$ as natural numbers and
        for every $d\in\naturals$ we have
            ($d$ divides $\naturalGcd{m}{n}$ as natural numbers iff
             $d$ divides $m$ as natural numbers and
             $d$ divides $n$ as natural numbers).
\end{proposition}
\begin{proof}
    Follows by
        \auto{gcd_good_pairs,times_tuple_intro,
            times_tuple_elim,pair_eq_iff}.
\end{proof}

\begin{proposition}\label{gcd_good_pair_zero_right}
    Let $m$ be a natural number.
    Then $(m,\emptyset)\in\naturalGcdGoodPairs$.
\end{proposition}
\begin{proof}
    $\emptyset$ is a natural number by \obv{omega_empty}.
    $\naturalGcd{m}{\emptyset}=m$
        by \obv{natural_gcd_zero_right}.
    $m$ divides $m$ as natural numbers
        by \obv{natural_divides_itself}.
    $m$ divides $\emptyset$ as natural numbers
        by \obv{every_natural_divides_zero}.
    Show that for every $d\in\naturals$ we have
        ($d$ divides $\naturalGcd{m}{\emptyset}$ as natural numbers iff
         $d$ divides $m$ as natural numbers and
         $d$ divides $\emptyset$ as natural numbers).
    \begin{subproof}
        Fix $d\in\naturals$.
        $d$ divides $\emptyset$ as natural numbers
            by \obv{every_natural_divides_zero}.
        Follows by \auto.
    \end{subproof}
    Follows by \auto{gcd_good_pair_iff}.
\end{proof}

\begin{corollary}\label{gcd_good_pair_zero_left}
    Let $n$ be a natural number.
    Then $(\emptyset,n)\in\naturalGcdGoodPairs$.
\end{corollary}
\begin{proof}
    $\emptyset$ is a natural number by \obv{omega_empty}.
    $\naturalGcd{\emptyset}{n}=n$
        by \obv{natural_gcd_zero_left}.
    $n$ divides $n$ as natural numbers
        by \obv{natural_divides_itself}.
    $n$ divides $\emptyset$ as natural numbers
        by \obv{every_natural_divides_zero}.
    Show that for every $d\in\naturals$ we have
        ($d$ divides $\naturalGcd{\emptyset}{n}$ as natural numbers iff
         $d$ divides $\emptyset$ as natural numbers and
         $d$ divides $n$ as natural numbers).
    \begin{subproof}
        Fix $d\in\naturals$.
        $d$ divides $\emptyset$ as natural numbers
            by \obv{every_natural_divides_zero}.
        Follows by \auto.
    \end{subproof}
    Follows by \auto{gcd_good_pair_iff}.
\end{proof}

\begin{proposition}\label{gcd_good_pair_symmetric}
    Let $m,n$ be natural numbers.
    Suppose $(m,n)\in\naturalGcdGoodPairs$.
    Then $(n,m)\in\naturalGcdGoodPairs$.
\end{proposition}
\begin{proof}
    $\naturalGcd{m}{n}=\naturalGcd{n}{m}$
        by \obv{natural_gcd_symmetric}.
    Follows by \auto{gcd_good_pair_iff}.
\end{proof}

\begin{theorem}\label{gcd_good_pair_sum_right}
    Let $m,k$ be natural numbers.
    Suppose $m\neq\emptyset$ and
        $(m,k)\in\naturalGcdGoodPairs$.
    Then $(m,m+k)\in\naturalGcdGoodPairs$.
\end{theorem}
\begin{proof}
    $m+k$ is a natural number by \obv{num_addition_on_naturals}.
    $\naturalGcd{m}{m+k}=\naturalGcd{m}{k}$
        by \obv{natural_gcd_subtract_sum}.
    $\naturalGcd{m}{k}$ divides $m$ as natural numbers and
        $\naturalGcd{m}{k}$ divides $k$ as natural numbers
        by \auto{gcd_good_pair_iff}.
    Hence $\naturalGcd{m}{m+k}$ divides $m$ as natural numbers.
    $\naturalGcd{m}{m+k}$ divides $k$ as natural numbers.
    $\naturalGcd{m}{m+k}$ divides $m+k$ as natural numbers
        by \obv{common_divisor_divides_sum}.
    Show that for every $d\in\naturals$ we have
        ($d$ divides $\naturalGcd{m}{m+k}$ as natural numbers iff
         $d$ divides $m$ as natural numbers and
         $d$ divides $m+k$ as natural numbers).
    \begin{subproof}
        Fix $d\in\naturals$.
        ($d$ divides $\naturalGcd{m}{k}$ as natural numbers iff
         $d$ divides $m$ as natural numbers and
         $d$ divides $k$ as natural numbers)
            by \auto{gcd_good_pair_iff}.
        ($d$ divides $m$ as natural numbers and
         $d$ divides $m+k$ as natural numbers)
            iff
        ($d$ divides $m$ as natural numbers and
         $d$ divides $k$ as natural numbers)
            by \obv{common_divisor_sum_iff_difference}.
    \end{subproof}
    Follows by \auto{gcd_good_pair_iff}.
\end{proof}

\begin{theorem}\label{every_natural_pair_is_gcd_good}
    Let $m,n$ be natural numbers.
    Then $(m,n)\in\naturalGcdGoodPairs$.
\end{theorem}
\begin{proof}
    Suppose not.
    Let $B=\{v\in\naturals\mid
        \text{there exists $u\in\naturals$ such that
            $(u,v)\notin\naturalGcdGoodPairs$}\}$.
    $n\in B$.
    Hence $B$ is inhabited.
    Take $q$ such that $q$ is an \in-minimal element of $B$
        by \obv{regularity}.
    Then $q\in B$ and $q\notmeets B$.
    Hence $q$ is a natural number and
        there exists $u\in\naturals$ such that
            $(u,q)\notin\naturalGcdGoodPairs$.

    Let $A=\{u\in\naturals\mid
        (u,q)\notin\naturalGcdGoodPairs\}$.
    $A$ is inhabited by \auto.
    Take $p$ such that $p$ is an \in-minimal element of $A$
        by \obv{regularity}.
    Then $p\in A$ and $p\notmeets A$.
    Hence $p$ is a natural number and
        $(p,q)\notin\naturalGcdGoodPairs$.

    Show that $q\neq\emptyset$.
    \begin{subproof}
        Suppose $q=\emptyset$.
        Then $(p,q)\in\naturalGcdGoodPairs$
            by \auto{gcd_good_pair_zero_right}.
        Contradiction.
    \end{subproof}
    Show that $p\neq\emptyset$.
    \begin{subproof}
        Suppose $p=\emptyset$.
        Then $(p,q)\in\naturalGcdGoodPairs$
            by \auto{gcd_good_pair_zero_left}.
        Contradiction.
    \end{subproof}
    $p$ and $q$ are ordinals by \obv{natural_number_is_ordinal}.

    \begin{byCase}
        \caseOf{$p\subseteq q$.}
            Take a natural number $k$ such that $q=p+k$
                by \obv{natural_additive_difference}.
            $k\in p+k$
                by \obv{positive_summand_makes_strict_sum}.
            Hence $k\in q$.
            Thus $k\notin B$ by \auto{disjoint}.
            Hence $(p,k)\in\naturalGcdGoodPairs$
                by \auto.
            Thus $(p,p+k)\in\naturalGcdGoodPairs$
                by \obv{gcd_good_pair_sum_right}.
            Hence $(p,q)\in\naturalGcdGoodPairs$.
            Contradiction.
        \caseOf{$\lnot(p\subseteq q)$.}
            Hence $q\subseteq p$
                by \obv{ordinal_subseteq_or_subseteq}.
            Take a natural number $k$ such that $p=q+k$
                by \obv{natural_additive_difference}.
            $k\in q+k$
                by \obv{positive_summand_makes_strict_sum}.
            Hence $k\in p$.
            Thus $k\notin A$ by \auto{disjoint}.
            Hence $(k,q)\in\naturalGcdGoodPairs$
                by \auto.
            Hence $(q,k)\in\naturalGcdGoodPairs$
                by \obv{gcd_good_pair_symmetric}.
            Thus $(q,q+k)\in\naturalGcdGoodPairs$
                by \obv{gcd_good_pair_sum_right}.
            Hence $(q,p)\in\naturalGcdGoodPairs$.
            Thus $(p,q)\in\naturalGcdGoodPairs$
                by \obv{gcd_good_pair_symmetric}.
            Contradiction.
    \end{byCase}
\end{proof}

\begin{theorem}\label{natural_gcd_common_divisor}
    Let $m,n$ be natural numbers.
    Then $\naturalGcd{m}{n}$ divides $m$ as natural numbers and
        $\naturalGcd{m}{n}$ divides $n$ as natural numbers.
\end{theorem}
\begin{proof}
    $(m,n)\in\naturalGcdGoodPairs$
        by \obv{every_natural_pair_is_gcd_good}.
    Follows by \auto{gcd_good_pair_iff}.
\end{proof}

\begin{theorem}\label{natural_gcd_universal}
    Let $d,m,n$ be natural numbers.
    Then $d$ divides $\naturalGcd{m}{n}$ as natural numbers iff
        $d$ divides $m$ as natural numbers and
        $d$ divides $n$ as natural numbers.
\end{theorem}
\begin{proof}
    $(m,n)\in\naturalGcdGoodPairs$
        by \obv{every_natural_pair_is_gcd_good}.
    Follows by \auto{gcd_good_pair_iff}.
\end{proof}

\begin{proposition}\label{natural_gcd_off_domain}
    Suppose $m$ is not a natural number or $n$ is not a natural number.
    Then $\naturalGcd{m}{n}=\emptyset$.
\end{proposition}
\begin{proof}
    Show that $\naturalCommonDivisors{m}{n}=\emptyset$.
    \begin{subproof}
        Show that every element of $\naturalCommonDivisors{m}{n}$
            is an element of $\emptyset$.
        \begin{subproof}
            Fix $d\in\naturalCommonDivisors{m}{n}$.
            Contradiction.
        \end{subproof}
        Hence $\naturalCommonDivisors{m}{n}\subseteq\emptyset$ by
            \obv{subseteq}.
        Follows by \obv{subseteq_emptyset_iff}.
    \end{subproof}
    Follows by \obv{natural_gcd,unions_emptyset}.
\end{proof}

\begin{proposition}\label{coprime_naturals_symmetric}
    Let $m,n$ be natural numbers.
    Suppose $m$ is coprime to $n$ as natural numbers.
    Then $n$ is coprime to $m$ as natural numbers.
\end{proposition}
\begin{proof}
    Follows by
        \auto{coprime_naturals,
            natural_gcd_symmetric}.
\end{proof}

\subsection{Prime divisors of products}

\begin{proposition}\label{divisor_divides_multiple}
    Let $d,a,b$ be natural numbers.
    Suppose $d$ divides $a$ as natural numbers.
    Then $d$ divides $\naturalMul{a}{b}$ as natural numbers.
\end{proposition}
\begin{proof}
    Take a natural number $c$ such that $a=\naturalMul{d}{c}$
        by \obv{natural_divisibility}.
    $\naturalMul{c}{b}$ is a natural number
        by \obv{natural_multiplication_closed}.
    $\naturalMul{a}{b}=\naturalMul{d}{\naturalMul{c}{b}}$
        by \auto{natural_multiplication_associative}.
    $\naturalMul{a}{b}$ is a natural number
        by \obv{natural_multiplication_closed}.
    Follows by \obv{natural_divisibility}.
\end{proof}

\begin{theorem}\label{natural_division_algorithm}
    Let $a,b$ be natural numbers.
    Suppose $b\neq\emptyset$.
    Then there exist $q,r\in\naturals$ such that
        $a=\naturalMul{b}{q}+r$ and $r\in b$.
\end{theorem}
\begin{proof}
    Let $C=\{n\in\naturals\mid
        \text{there exist $q,r\in\naturals$ such that
            $n=\naturalMul{b}{q}+r$ and $r\in b$}\}$.
    Show that $\emptyset\in C$.
    \begin{subproof}
        $\emptyset\in b$ by \obv{natural_nonzero_contains_zero}.
        $\naturalMul{b}{\emptyset}=\emptyset$
            by \obv{natural_multiplication_zero_right}.
        $\emptyset+\emptyset=\emptyset$
            by \obv{num_addition_zero_right,omega_empty}.
        Follows by \auto{omega_empty}.
    \end{subproof}
    Show that for every $n\in C$ we have $\suc{n}\in C$.
    \begin{subproof}
        Fix $n\in C$.
        Take $q,r\in\naturals$ such that
            $n=\naturalMul{b}{q}+r$ and $r\in b$.
        $\naturalMul{b}{q}$ is a natural number
            by \obv{natural_multiplication_closed}.
        $\suc{n}=\naturalMul{b}{q}+\suc{r}$
            by \auto{num_addition_successor_right}.
        $r$ and $b$ are ordinals by \obv{natural_number_is_ordinal}.
        Hence $\suc{r}\subseteq b$ by \obv{ordinal_suc_subseteq}.
        $\suc{r}$ is a natural number by \obv{omega_successor}.
        $\suc{r}$ is an ordinal by \obv{natural_number_is_ordinal}.
        \begin{byCase}
            \caseOf{$\suc{r}=b$.}
                $\suc{q}$ is a natural number by \obv{omega_successor}.
                $\naturalMul{b}{\suc{q}}
                    =\naturalMul{b}{q}+b$
                    by \obv{natural_multiplication_successor_right}.
                $\naturalMul{b}{\suc{q}}$ is a natural number
                    by \obv{natural_multiplication_closed}.
                $\naturalMul{b}{\suc{q}}+\emptyset
                    =\naturalMul{b}{\suc{q}}$
                    by \obv{num_addition_zero_right}.
                \begin{align*}
                    \suc{n}
                        &= \naturalMul{b}{q}+\suc{r}
                        \\
                        &= \naturalMul{b}{q}+b
                        \\
                        &= \naturalMul{b}{\suc{q}}
                        \\
                        &= \naturalMul{b}{\suc{q}}+\emptyset
                \end{align*}
                $\emptyset\in b$
                    by \obv{natural_nonzero_contains_zero}.
                Follows by \auto{omega_empty}.
            \caseOf{$\suc{r}\neq b$.}
                Hence $\suc{r}\in b$
                    by \obv{ordinal_subseteq_implies_elem_or_eq}.
                Follows by \auto{omega_successor}.
        \end{byCase}
    \end{subproof}
    Hence $C$ is an inductive set by \obv{inductive_set}.
    Hence every natural number is an element of $C$
        by \obv{omega_minimal}.
    Thus $a\in C$.
\end{proof}

\begin{definition}\label{positive_multiplier_set}
    $\naturalPositiveMultipliers{d}{a}=
        \{t\in\naturals\mid
            \text{$t\neq\emptyset$ and
                $d$ divides $\naturalMul{a}{t}$ as natural numbers}\}$.
\end{definition}

\begin{proposition}\label{positive_multiplier_iff}
    $t\in\naturalPositiveMultipliers{d}{a}$ iff
        $t$ is a natural number and $t\neq\emptyset$ and
        $d$ divides $\naturalMul{a}{t}$ as natural numbers.
\end{proposition}
\begin{proof}
    Follows by \auto{positive_multiplier_set}.
\end{proof}

\begin{proposition}\label{positive_multiplier_intro}
    Let $t$ be a natural number.
    Suppose $t\neq\emptyset$ and
        $d$ divides $\naturalMul{a}{t}$ as natural numbers.
    Then $t\in\naturalPositiveMultipliers{d}{a}$.
\end{proposition}
\begin{proof}
    Follows by \obv{positive_multiplier_iff}.
\end{proof}

\begin{theorem}\label{least_positive_multiplier_divides_all}
    Let $d,a,r$ be natural numbers.
    Suppose $r$ is an \in-minimal element of
        $\naturalPositiveMultipliers{d}{a}$.
    Then for every $t\in\naturalPositiveMultipliers{d}{a}$ we have
        $r$ divides $t$ as natural numbers.
\end{theorem}
\begin{proof}
    $r\in\naturalPositiveMultipliers{d}{a}$ and
        $r\notmeets\naturalPositiveMultipliers{d}{a}$.
    Hence $r\neq\emptyset$ and
        $d$ divides $\naturalMul{a}{r}$ as natural numbers
        by \obv{positive_multiplier_iff}.
    Fix $t\in\naturalPositiveMultipliers{d}{a}$.
    Then $t$ is a natural number and $t\neq\emptyset$ and
        $d$ divides $\naturalMul{a}{t}$ as natural numbers
        by \obv{positive_multiplier_iff}.
    Take $q,s\in\naturals$ such that
        $t=\naturalMul{r}{q}+s$ and $s\in r$
        by \obv{natural_division_algorithm}.

    $d$ divides $\naturalMul{\naturalMul{a}{r}}{q}$ as natural numbers
        by \obv{divisor_divides_multiple,
            natural_multiplication_closed}.
    $\naturalMul{\naturalMul{a}{r}}{q}
        =\naturalMul{a}{\naturalMul{r}{q}}$
        by \obv{natural_multiplication_associative}.
    Hence $d$ divides $\naturalMul{a}{\naturalMul{r}{q}}$
        as natural numbers.
    $\naturalMul{r}{q}$ is a natural number
        by \obv{natural_multiplication_closed}.
    $\naturalMul{a}{t}
        =\naturalMul{a}{\naturalMul{r}{q}+s}$.
    $\naturalMul{a}{\naturalMul{r}{q}+s}
        =\naturalMul{a}{\naturalMul{r}{q}}+\naturalMul{a}{s}$
        by \obv{natural_multiplication_distributive_left}.
    $\naturalMul{a}{t}
        =\naturalMul{a}{\naturalMul{r}{q}}+\naturalMul{a}{s}$.
    Hence $d$ divides $\naturalMul{a}{s}$ as natural numbers
        by \auto{common_divisor_of_sum_difference,
            natural_multiplication_closed}.

    Show that $s=\emptyset$.
    \begin{subproof}
        Suppose not.
        Hence $s\in\naturalPositiveMultipliers{d}{a}$
            by \obv{positive_multiplier_intro}.
        Contradiction by \auto{disjoint}.
    \end{subproof}
    Hence $t=\naturalMul{r}{q}$
        by \auto{num_addition_zero_right}.
    Follows by \auto{natural_divisibility}.
\end{proof}

\begin{theorem}\label{prime_divides_natural_product}
    Let $p,a,b$ be natural numbers.
    Suppose $p$ is a prime natural number and
        $p$ divides $\naturalMul{a}{b}$ as natural numbers.
    Then $p$ divides $a$ as natural numbers or
        $p$ divides $b$ as natural numbers.
\end{theorem}
\begin{proof}
    \begin{byCase}
        \caseOf{$p$ divides $a$ as natural numbers.}
            Follows.
        \caseOf{$p$ does not divide $a$ as natural numbers.}
            \begin{byCase}
                \caseOf{$b=\emptyset$.}
                    Hence $p$ divides $b$ as natural numbers
                        by \obv{every_natural_divides_zero}.
                \caseOf{$b\neq\emptyset$.}
                    Show that $p\neq\emptyset$.
                    \begin{subproof}
                        Suppose not.
                        $\naturalOne\in p$
                            by \obv{prime_natural}.
                        Contradiction by \obv{emptyset}.
                    \end{subproof}

                    Let $S=\naturalPositiveMultipliers{p}{a}$.
                    $p$ divides $p$ as natural numbers
                        by \obv{natural_divides_itself}.
                    Hence $p$ divides $\naturalMul{p}{a}$
                        as natural numbers
                        by \obv{divisor_divides_multiple}.
                    $\naturalMul{p}{a}=\naturalMul{a}{p}$
                        by \obv{natural_multiplication_commutative}.
                    Hence $p$ divides $\naturalMul{a}{p}$
                        as natural numbers.
                    $p\in\naturalPositiveMultipliers{p}{a}$ by
                        \obv{positive_multiplier_intro}.
                    Hence $p\in S$.
                    Hence $S$ is inhabited.
                    $b\in\naturalPositiveMultipliers{p}{a}$ by
                        \obv{positive_multiplier_intro}.
                    Hence $b\in S$.

                    Take $r$ such that $r$ is an \in-minimal element of $S$
                        by \obv{regularity}.
                    $r\in S$ and $r\notmeets S$.
                    Hence $r$ is a natural number and $r\neq\emptyset$
                        by \auto{positive_multiplier_iff}.
                    $r$ divides $p$ as natural numbers and
                        $r$ divides $b$ as natural numbers
                        by \obv{least_positive_multiplier_divides_all}.
                    Hence $r=\naturalOne$ or $r=p$
                        by \obv{prime_natural}.

                    \begin{byCase}
                        \caseOf{$r=\naturalOne$.}
                            $p$ divides $\naturalMul{a}{r}$
                                as natural numbers
                                by \auto{positive_multiplier_iff}.
                            $\naturalMul{a}{r}=a$ by
                                \auto{natural_one,
                                    natural_multiplication_one_right}.
                            Hence $p$ divides $a$ as natural numbers.
                            Contradiction.
                        \caseOf{$r\neq\naturalOne$.}
                            Hence $r=p$.
                            Thus $p$ divides $b$ as natural numbers.
                    \end{byCase}
            \end{byCase}
    \end{byCase}
\end{proof}
